\documentclass[11pt]{article}

% Change "review" to "final" to generate the final (sometimes called camera‑ready)
% version. Change to "preprint" to generate a non‑anonymous version with page numbers.
\usepackage[review]{acl}

% Standard package includes
\usepackage{times}
\usepackage{latexsym}

% For proper rendering and hyphenation of words containing Latin characters (including
% in bib files). See https://www.latex-project.org/help/documentation/encguide.pdf for
% other character sets
\usepackage[T1]{fontenc}

% This assumes your files are encoded as UTF8
\usepackage[utf8]{inputenc}

% This is not strictly necessary, and may be commented out, but it will improve the
% layout of the manuscript, and will typically save some space.
\usepackage{microtype}

% This is also not strictly necessary, and may be commented out. However, it will
% improve the aesthetics of text in the typewriter font.
\usepackage{inconsolata}

%Including images in your LaTeX document requires adding additional package(s)
\usepackage{graphicx}


\title{MultiIFEval: An Instruction Following Benchmark in 305 Languages}

\author{Dan Saattrup Smart \\
  Alexandra Institute \\
  \texttt{dan.smart@alexandra.dk}\\
}

\begin{document}
\maketitle


\begin{abstract}

We introduce MultiIFEval, a large‑scale instruction‑following benchmark that spans 305
languages. Starting from the English IFEval dataset \cite{zhou2023ifeval}, we translate
and localise each prompt–response pair into 305 target languages using an LLM-based
pipeline with Wikipedia-grounded localisation. We create a manually translated Danish
variant of the dataset and conduct human expert validation to assess quality differences
between machine-translated and manually-translated Danish. We evaluate models on both
machine-translated and manually-translated Danish, as well as on Estonian \cite{dorkin2025ifeval-et}
and French \cite{alhajar2024ifeval-fr} variants from prior work. Finally, we evaluate
multiple models across all 305 languages to establish a comprehensive baseline for
cross-lingual instruction following capability.

\end{abstract}


\section{Introduction}
\label{sec:intro}

Instruction following evaluation has become central to developing practical AI systems
capable of understanding and executing natural language commands. The release of IFEval
\cite{zhou2023ifeval} established a standardised framework for measuring how well models
can follow verifiable instructions in English, providing a rigorous benchmark that goes
beyond traditional task completion metrics. However, IFEval's English-only scope reveals
a critical limitation: as large language models are increasingly deployed across diverse
linguistic contexts, evaluation benchmarks confined to English fail to capture
performance disparities that arise from multilingual instruction following.

The gap between monolingual and multilingual evaluation has only partially been addressed
by recent work. M-IFEval \cite{li2025xifbench} expanded to 3 languages, XIFBench
\cite{he2024multi} covered 6 languages, and Multi-IF \cite{he2024multi} reached 8
languages. More recently, Marco-Bench-MIF \cite{zeng-etal-2025-marco} extended coverage
to 30 languages. While these efforts represent important progress, they collectively
cover fewer than 1\% of the world's approximately 7,000 languages, leaving the vast
majority of linguistic communities without instruction-following evaluation resources.
This limitation is particularly acute for low-resource languages, where model behaviour
may diverge significantly from high-resource settings due to differences in training
data availability, tokenisation quality, and linguistic structure.

To bridge this gap, we introduce \textbf{MultiIFEval}, a benchmark that expands IFEval
to 305 target languages, representing an order-of-magnitude increase in multilingual
coverage. Our methodology leverages large language models to both translate and localise
every English prompt–response pair from the original IFEval dataset into each target
language. Crucially, our localisation pipeline grounds translations in Wikipedia
articles in the target language, ensuring that cultural references, proper nouns, and
contextual examples are adapted appropriately rather than translated literally. The
resulting dataset spans high-resource languages such as Spanish, Mandarin, and Arabic,
as well as low-resource languages including Xhosa, Quechua, and Samoan, providing an
unprecedented testbed for evaluating multilingual instruction following at scale.

In addition to machine-translated versions, we provide a manually translated and
localised Danish subset curated by native speakers and validated by domain experts. This
dual approach enables systematic comparison between machine-translated and
human-translated instruction-following benchmarks, allowing us to assess whether
translation quality meaningfully impacts model evaluation outcomes—an aspect largely
unexplored in prior multilingual benchmark work. We also incorporate manually translated
Estonian \cite{dorkin2025ifeval-et} and French \cite{alhajar2024ifeval-fr} variants
from recent independent efforts, enabling cross-lingual analysis across multiple
European languages with human-validated translations.

Our main contributions are methodological:

\begin{itemize}
\item We present a 305-language instruction-following benchmark, expanding multilingual
  coverage from the prior maximum of 30 languages to 305, spanning languages across all
  inhabited continents and representing diverse language families and resource levels.

\item We create a manually translated and localised Danish variant, enabling direct
  comparison between machine translation and human translation quality for instruction
  following evaluation, with expert validation to assess the validity of
  LLM-based translation pipelines.

\item We develop a methodology for LLM-based translation with Wikipedia-grounded
  localisation, providing a reproducible approach for adapting instruction-following
  benchmarks to new languages while preserving cultural and contextual appropriateness.

\item We establish an evaluation framework for cross-lingual instruction following that
  will be applied across all 305 languages, enabling systematic analysis of how model
  performance varies across language families, resource levels, and translation methods.
\end{itemize}


\section{Dataset Generation}
\label{sec:dataset-generation}

The MultiIFEval dataset is generated via prompting a large language model in a zero-shot
fashion, for each sample of IFEval. The prompt contains both the instruction to
translate and localise each English sample. To reduce the risk of hallucination when
localising the examples, as well as to increase the quality of the translations, we
include for each IFEval sample a randomly chosen Wikipedia page in the target language.
We use structured generation to ensure that the output is of the correct format. The
full prompt can be found in Figure~\ref{fig:prompt}.

The structure of the generated \texttt{kwargs} deviate from the input \texttt{kwargs},
and the reason for this is that the language model we used for generation did not
support dictionaries as values in the JSON schema used in structured generation, so the
is merely a rephrasing of the keyword arguments to obey this constraint.

TODO


\begin{figure*}[h]
\centering
\small
\begin{verbatim}
You are a professional translator from English to {target_language} (language code:
'{target_language_code}').

Here is an instruction-following example in English:

<example>
{"prompt": <prompt>, "instruction_id_list": <instruction_id_list>, "kwargs": <kwargs>}
</example>

You need to translate the example to {target_language}. This means the following:

1. The `prompt` should be translated to {target_language}. You should localise these to
   {target_language} and its country's culture.
2. The `instruction_id_list` should remain exactly the same.
3. The keys in the kwargs dictionaries should be completely unchanged, and the values
   should _only_ be changed if they contain English words or phrases, which would need
   to be changed to the equivalent in {target_language}.

Here is an example of some text written in {language.name}:

<{target_language_code}_example>
{wikipedia_article_in_target_language}
</{target_language_code}_example>

You should return the translated example in JSON format, with the following structure:

- `new_prompt` (str): The translated prompt.
- `new_instruction_id_list` (list[str]): The translated instruction ID list.
- `new_kwargs` (list[list[dict]]): The translated keyword arguments.

Here the `new_kwargs` must be a list of the same length as the
`new_instruction_id_list`, containing lists of dicts, each with keys `name` and `value`,
where `name` is the name of the keyword argument and `value` is the value of the keyword
argument. Each list of kwarg dicts should be the keyword arguments for the corresponding
instruction ID in the `new_instruction_id_list`.
\end{verbatim}
\caption{Prompt used to generate the MultiIFEval dataset.}
\label{fig:prompt}
\end{figure*}


\subsection{Manual Danish Translation}
\label{sec:manual-danish-translation}

Missing

\subsection{Machine Translated Dataset}
\label{sec:mt-dataset}

Missing


\section{Evaluations of Language Models}
\label{sec:evaluation}

Missing


\section{Comparison of Machine Translated and Manually Translated Datasets}
\label{sec:comparison}

Missing


\section{Conclusion}
\label{sec:conclusion}

Missing


\section{Acknowledgements}
\label{sec:acknowledgements}

Missing


\bibliography{custom,anthology}

\appendix

\section{Example Appendix}
\label{sec:appendix}

Missing

\end{document}
